\chapter{Introduction}

Introductory lines...

\section{Section-1 Name}

Some text here.

\begin{definition}
Some definition...
\end{definition}

\begin{theorem}\label{thm1}
Some theorem...
\end{theorem}

\begin{proof}
Proof is as follows...
\end{proof}

\begin{corollary}
A corollary to \autoref{thm1} is...
\end{corollary}

\begin{remark}
Some remark...
\end{remark}

\subsection{Equations and Math Examples}

Equations can be typed as follows:

\begin{equation}\label{eqn1}
f(x) = \frac{x^2 - 5x + 6}{(e^x - 2)/10} = 10 \times \frac{(x-2)(x-3)}{e^x-2}
\end{equation}

Referencing labelled objects: \autoref{eqn1}, or \autoref{thm1}.  

For multiline equations,
\begin{equation}\label{eqn2}
\text{Array in Math Mode} \quad
\left\{
\begin{array}{rcll}
-\Delta u + \lambda u &=& |u|^{p-2}, & \text{in } \Omega\\
u &\geq& 0, & u \in H_0^1(\Omega)
\end{array}
\right.
\end{equation}

% r, c, and l stand for the alignment operators for each column (right, centre, left), defined with '&' acting as the delimiter/column separators.

Using \verb+array+ in math mode or \verb+eqnarray+ is a quick and easy way to get the most customisable equation output, but is outdated and longer equations are prone to errors. Use of alternate multiline equation environments like \verb+multiline+(*), \verb+align+(*), \verb+gather+(*) or \verb+split+ in any math-mode environment is recommended.

\begin{align}\label{eq3}
g(\theta) &= \, i\theta &=& \, (i\theta) \times \ln e \nonumber \\
&= \, \ln(e^{i\theta}) &=& \, \ln(\cos\theta + i\sin\theta)
\end{align}

% \nonumber will suppress calling for an equation reference number
% Alternatively, starred environments (for example, align*) achieve the same result

\section{Section-2 Name}

Matrices in {\LaTeX} look like:

\begin{align*}
\begin{pmatrix}
\sin\theta & \cos\theta\\
-\cos\theta & \sin\theta
\end{pmatrix}
\times
\begin{pmatrix}
\sin\theta & \cos\theta\\
-\cos\theta & \sin\theta
\end{pmatrix} &=
\begin{pmatrix}
\sin^2\theta - \cos^2\theta & 2\cos\theta\sin\theta\\
-2\cos\theta\sin\theta & -\cos^2\theta + \sin^2\theta
\end{pmatrix} \\
&=
\begin{pmatrix}
-\cos2\theta & \sin2\theta\\
-\sin2\theta & -\cos2\theta
\end{pmatrix}
\end{align*}
\begin{multicols}{2}
[The brackets of a given matrix depend on the type of matrix called.]
\begin{flushright}
    Here is a quick truth table:
\end{flushright}
\columnbreak
\begin{tabular}{|c c c | c|}
\hline % Horizontal bar across table
$P$ & $Q$ & $\neg P$ & $\neg P \to (P \vee Q)$ \\ \hline
T & T & F & T \\
T & F & F & T \\
F & T & T & T \\
F & F & T & F \\ \hline
\end{tabular}
\end{multicols}

\begin{remark}
Defining a table like this does not count in the LoT; use the \verb+table+ environment instead.
\end{remark}

\begin{remark}
You can cite sources in footnotes as so\autocite{golub}.  
Ensure \verb+ref.bib+ is configured for biblatex. Disable \texttt{verbose} style to switch to inline references.
\end{remark}

\newpage % to fix the footnote getting "stuck" to the text above.

\subsection{Subsections}

\subsubsection{Subsubsection Example}

Subsubsections do not appear in the ToC and lack numbering\footnote{Footnotes work for comments as well. For more, see \url{https://www.overleaf.com/learn/latex/Footnotes}}.
To skip numbering in sections/subsections, use \verb+\section*+\{\textit{section\_name}\}.

\begin{theorem}\label{Th123}
Some theorem...
\end{theorem}

\begin{proof}
The proof is as follows...
\end{proof}

\begin{figure}[htbp]
\centering
\captionsetup{justification=centering}
\input{Images/Figures/3D_Cone}
\caption{3D Cone designed by Gene R. using TikZ, see \texttt{Images/Figures/3D\_Cone.tex}. Delete to save compile time.}
\end{figure}

\begin{remark}
Figures float by default. Position may differ from the order in the code. Use optional arguments [!htbp] (here, top, bottom, next page) to influence placement.
\end{remark}

To make your own commutative diagrams, consider using tools such as \href{graphviz.org}{GraphViz}, \href{q.uiver.app}{Quiver} and \href{ipe.otfried.org}{IPE}.

\newpage % Avoids footer sticking to text

%\subsection{Chemical Notations} % Uncomment 'CHEMICAL NOTATION' packages in Preamble
%\ce{HC#CH + H2O ->[Hg^++][18\% H2SO4, $90^\circ$] CH3 - CHO} \\
%\ce{H2 (g) + I2 (g) <=>[Formation][Decomposition] 2HI (g)}\\
%\begin{figure}[htbp]
%\centering
%\chemfig{*6((=O)-N(-H)-(-R_2)-(=O)-N(-H)-(-R_1)-)}
%\caption{Diketopiperazine}
%\end{figure}

\section{Sample Question and Proof}

Suppose $A_i$ is a connected subset of a topological space $X$ for $i=1,\ldots,n$, and $A_i \cap A_{i+1} \neq \emptyset$ for all $i$. 
Prove that $A = \bigcup_{i=1}^n A_i$ is connected.

\begin{proof}[Proof by Contradiction]
Assume $A$ is disconnected. $A$ can then be written as a union of two non-empty, disjoint, relatively open subsets, say, $X$ and $Y$. Take some $x \in X$ and some $y \in Y$, with $x \in A_j$ and $y \in A_k$ for some $j \le k$. Then
\begin{align}\label{sampleeq}
A_l \cap A_{l+1} \neq \emptyset \qquad & \forall \, l \in \{j,\ldots,k-1\} \nonumber\\
\Rightarrow \qquad \therefore \quad \bigcup_{i=j}^l A_i \text{ is connected } \quad & \forall \, l \in \{j,\ldots,k\}
\end{align}
Hence, $\bigcup_{i=1}^n A_i$ contains both $x$ and $y$ and is connected, contradicting our original assumption of the disjointness of $X$ and $Y$. Therefore, $A=\bigcup_{i=j}^k A_i$ is connected.
\end{proof}

\begin{remark}
\verb+\quad+, \verb+\qquad+, \verb+\,+ and \verb+\!+ are effective in adjusting spacing as needed.
\end{remark}
